
\documentclass{article}
\usepackage{colortbl}
\usepackage{makecell}
\usepackage{multirow}
\usepackage{supertabular}

\begin{document}

\newcounter{utterance}

\centering \large Interaction Transcript for game `cladder', experiment `full\_v1.5\_default', episode 7410 with Qwen3.6{-}27B.
\vspace{24pt}

{ \footnotesize  \setcounter{utterance}{1}
\setlength{\tabcolsep}{0pt}
\begin{supertabular}{c@{$\;$}|p{.15\linewidth}@{}p{.15\linewidth}p{.15\linewidth}p{.15\linewidth}p{.15\linewidth}p{.15\linewidth}}
    \# & \multicolumn{2}{c}{Player} && \multicolumn{2}{c}{Game Master} \\
    \hline

    \theutterance \stepcounter{utterance}  
    & & & \multicolumn{4}{p{0.6\linewidth}}{
        \cellcolor[rgb]{0.9,0.9,0.9}{
            \makecell[{{p{\linewidth}}}]{
                \texttt{\tiny{[P1$\langle$GM]}}
                \texttt{You are an expert in causal inference. The following question is not a typical commonsense query, but rather a meticulously designed question created by a professor specializing in causal inference, intended to assess the students' mastery of the course content.} \\
\texttt{Imagine a self{-}contained, hypothetical world with only the following conditions, and without any unmentioned factors or causal relationships: Wibl has a direct effect on zuph and uvzi. Zuph has a direct effect on uvzi. Method 1: We look at how zuph correlates with uvzi case by case according to wibl. Method 2: We look directly at how zuph correlates with uvzi in general. To understand how zuph affects uvzi, is it more correct to use the Method 1 than Method 2?} \\
\texttt{Start your answer with 'Yes' or 'No', followed by additional reasoning or evidence to support your explanation.} \\
            }
        }
    }
     \\ \\

    \theutterance \stepcounter{utterance}  
    & \multicolumn{4}{p{0.6\linewidth}}{
        \cellcolor[rgb]{0.561,0.737,0.561}{
            \makecell[{{p{\linewidth}}}]{
                \texttt{\tiny{[P1$\rangle$GM]}}
                \texttt{Yes} \\
\\ 
\texttt{**Reasoning:**} \\
\\ 
\texttt{To determine the correct method for estimating the causal effect of **Zuph** on **Uvzi**, we must analyze the causal structure provided:} \\
\\ 
\texttt{1.  **Causal Structure:**} \\
\texttt{    *   \$Wibl \textbackslash{}rightarrow Zuph\$} \\
\texttt{    *   \$Wibl \textbackslash{}rightarrow Uvzi\$} \\
\texttt{    *   \$Zuph \textbackslash{}rightarrow Uvzi\$} \\
\\ 
\texttt{    In this structure, **Wibl** is a **common cause** (confounder) of both the treatment (\$Zuph\$) and the outcome (\$Uvzi\$).} \\
\\ 
\texttt{2.  **Analysis of Method 2 (Marginal Correlation):**} \\
\texttt{    *   Method 2 looks at the direct correlation between \$Zuph\$ and \$Uvzi\$ in general (i.e., without conditioning on \$Wibl\$).} \\
\texttt{    *   Because \$Wibl\$ affects both \$Zuph\$ and \$Uvzi\$, there is a non{-}causal backdoor path: \$Zuph \textbackslash{}leftarrow Wibl \textbackslash{}rightarrow Uvzi\$.} \\
\texttt{    *   This backdoor path creates a spurious association between \$Zuph\$ and \$Uvzi\$. Therefore, the marginal correlation (Method 2) will be biased by the confounding effect of \$Wibl\$. It does not isolate the direct causal effect of \$Zuph\$ on \$Uvzi\$.} \\
\\ 
\texttt{3.  **Analysis} \\
            }
        }
    }
    & & \\ \\

    \theutterance \stepcounter{utterance}  
    & & & \multicolumn{2}{p{0.3\linewidth}}{
        \cellcolor[rgb]{0.9,0.9,0.9}{
            \makecell[{{p{\linewidth}}}]{
                \texttt{\tiny{[GM$|$GM]}}
                \texttt{Yes} \\
            }
        }
    }
    & & \\ \\

    \theutterance \stepcounter{utterance}  
    & & & \multicolumn{2}{p{0.3\linewidth}}{
        \cellcolor[rgb]{0.9,0.9,0.9}{
            \makecell[{{p{\linewidth}}}]{
                \texttt{\tiny{[GM$|$GM]}}
                \texttt{yes} \\
            }
        }
    }
    & & \\ \\

    \theutterance \stepcounter{utterance}  
    & & & \multicolumn{2}{p{0.3\linewidth}}{
        \cellcolor[rgb]{0.9,0.9,0.9}{
            \makecell[{{p{\linewidth}}}]{
                \texttt{\tiny{[GM$|$GM]}}
                \texttt{game\_result = WIN} \\
            }
        }
    }
    & & \\ \\

\end{supertabular}
}

\end{document}
